\chapter{Live Experiments}\label{ch:rs:live-experiments}

A new execution algorithm handles ten per cent of the firm's orders for two weeks, 2\,485 of 24\,000, and saves 0.55 basis point of \omterm{def:rs:simulation-versus-live:calibration}{implementation shortfall}
on average, with a standard error of 0.48. The 95\% interval runs from $-1.49$ to 0.39: the experiment cannot tell a saving from a cost. At half the flow it
would need 77 trading days to detect a saving of 0.3 basis point, and 38 with the variance reduction of this chapter. Even a perfect answer would be to a
narrower question than the desk's: in the chapter's simulated flow the new algorithm saves 0.30 basis point on its own orders but adds 0.25 to the firm's
other orders in the same stock, side and day, so that switching everything over saves 0.08. Research reaches production through \omterm{def:rs:live-experiments:paper}{paper trading}, canaries
and \omterm{def:rs:live-experiments:canary}{staged rollouts}, and is then measured on live flow; this chapter covers both, with \texttt{firm.abtest}.

\section{Paper trading and shadow mode}

\begin{definition}[Paper trading]\label{def:rs:live-experiments:paper}
\emph{Paper trading}\index{paper trading} runs a new strategy in real time on live market data: it computes its signals and generates its orders, which are
recorded and checked but never sent, and fills them with a \omterm{def:rs:event-driven-backtests:fill}{fill model}.
\end{definition}

\omterm{def:rs:live-experiments:paper}{Paper trading} tests everything that does not depend on the market's reaction. The live feeds arrive with their real timing and gaps; the signals computed
live can be compared with the same signals recomputed by the research code on the same data, value by value; the orders pass or fail the pre-trade risk
checks; the strategy lives through the operational calendar (start-up and shutdown, corporate actions, holidays, the index rebalance) that a \omterm{def:rs:vectorised-backtests:backtest}{backtest}
compresses into rows. What \omterm{def:rs:live-experiments:paper}{paper trading} cannot test is the fills. It is a simulator running on live data, with the blind spots of every simulator: chapter~19
showed that a replay without the strategy's own orders misses the trades that stopped at them and the levels where they were alone. \omterm{def:rs:simulation-versus-live:reconciliation}{Shadow trading}
(chapter~19) compares the simulator with the live trading of an existing strategy; \omterm{def:rs:live-experiments:paper}{paper trading} is the stage before a new strategy has any.

The exit condition of \omterm{def:rs:live-experiments:paper}{paper trading} is therefore not its P\&L. It is parity: live signals equal to the research signals within a stated tolerance, every
order it would have sent accepted by the risk checks, and a full operational cycle crossed without manual repair. A strategy whose paper P\&L is good and
whose live signals differ from research has not been tested; it has been replaced by another one.

\section{Canaries and staged rollouts}

\begin{definition}[Canary deployment, staged rollout]\label{def:rs:live-experiments:canary}
A \emph{canary deployment}\index{canary deployment} runs a new version of a trading component on a small share of the production flow (the canary) alongside
the old version, and compares the two before the new one receives more. A \emph{staged rollout}\index{staged rollout} expands the share in steps (for instance
1\%, 10\%, 50\%, 100\%), each step gated by stated checks, with a tested path back to the old version at every step.
\end{definition}

The most expensive deployment failure in the public record was a staged deployment. From 27 July 2012 Knight Capital placed new order-routing code on a
limited number of servers on successive days, ahead of a new NYSE programme starting on 1 August. A technician did not copy the new code to one of the eight
servers; no second person reviewed the deployment, and no written procedure required it. On 1 August the eighth server ran an old, defective function that
the new code's repurposed flag reactivated: for 212 incoming parent orders it sent millions of child orders, 4 million executions in 154 stocks for more than
397 million shares in about 45 minutes. Knight was left with a net long position of about \$3.5 billion in 80 stocks and a net short of about \$3.15 billion in
74, and realised a loss of \$460 million. Before the open, its systems had sent 97 automated e-mails citing the old function by name; they were not designed as
alerts and nobody read them. The SEC's order of October 2013 imposed a \$12 million penalty.

Three lessons follow, none of them statistical. A stage is only as good as its check of the deployed state: every server, process and configuration runs the
version the stage says it runs. An alert is a message designed and routed to be acted on, not a log line that happens to be e-mailed. And a stage must be capped
in what it can do: a canary with 1\% of the flow must hold 1\% of the risk limits, and invariants such as ``child quantity never exceeds the parent's'' must be
checked on every order, where they fire on the first violation. The statistical comparison of canary and control is for a different class of defect.

\begin{proposition}[How long a statistical canary takes]\label{prop:rs:live-experiments:canary}
A canary receives a share $s$ of a flow; a defect adds $\delta$ to a per-order metric of standard deviation $\sigma$; the alarm is a one-sided $z$-test of the
canary's mean against the control's at threshold $z$. The expected alarm comes after about
\[
n \approx \frac{z^2\sigma^2}{\delta^2(1 - s)}
\]
canary orders: nearly independent of the share when it is small.
\end{proposition}

\begin{proof}
After $n$ canary orders the control has $m = n(1-s)/s$, and the difference of means has variance $\sigma^2/n + \sigma^2/m = \sigma^2/(n(1-s))$. The expected
statistic $\delta\sqrt{n(1-s)}/\sigma$ reaches $z$ at the stated $n$.
\end{proof}

With the chapter's flow ($\sigma = 23.1$ basis points, 2\,400 orders a day), a defect costing 5 basis points an order and an alarm at $z = 3$, a 1\% canary
raises the alarm after 194 damaged orders and 8.1 days, a 10\% canary after 214 orders and 0.89 day, a 50\% canary after 385 orders and 0.32 day. A smaller
canary does not reduce the damage done before a statistical alarm; it spreads it over more time. What a small canary buys is a small exposure for defects
caught by invariants and limits, and time for people to look.

\begin{omfigure}
\begin{tikzpicture}[stage/.style={draw=omDef,fill=omDef!8,rounded corners=2pt,text width=1.85cm,align=center,font=\scriptsize,minimum height=1.1cm},
  gate/.style={text width=1.95cm,align=center,font=\scriptsize,text=black!70},arr/.style={-{Stealth},thick,omDef}]
\node[stage] (r) at (0,0) {research\\backtests};
\node[stage] (p) at (2.35,0) {paper trading\\(no orders)};
\node[stage] (c) at (4.7,0) {canary\\1\% of flow};
\node[stage] (a) at (7.05,0) {A/B test\\by stock-day};
\node[stage] (h) at (9.4,0) {50\%};
\node[stage] (f) at (11.75,0) {full rollout};
\foreach \x/\y in {r/p,p/c,c/a,a/h,h/f} \draw[arr] (\x) -- (\y);
\node[gate,below=3pt of r] {PBO, deflated SR, cost model (ch.~20, 27)};
\node[gate,below=3pt of p] {signal parity, risk checks, a full calendar};
\node[gate,below=3pt of c] {deployed state verified, invariants, 1\% of limits};
\node[gate,below=3pt of a] {effect and power, guardrail metrics};
\node[gate,below=3pt of h] {monitoring, capacity (ch.~28)};
\node[gate,below=3pt of f] {kill switch, path back tested};
\draw[omThm,thick,-{Stealth},dashed] (f.north) to[bend right=12] node[above,font=\scriptsize,text=omThm] {roll back at any stage} (c.north);
\end{tikzpicture}

\omcaption{From research to production: each stage is gated by checks it can actually perform. \omterm{def:rs:live-experiments:paper}{Paper trading} cannot measure fills; a canary's statistics
cannot catch a runaway order in time, its invariants and limits can.}\label{fig:rs:live-experiments:stages}
\end{omfigure}

\section{Randomised A/B tests on live flow}

\begin{definition}[A/B test, randomisation unit, interference, sample ratio mismatch]\label{def:rs:live-experiments:ab}
An \emph{A/B test}\index{A/B test} compares two versions of a component by assigning units of live flow to them at random and comparing an outcome. The
\emph{randomisation unit}\index{randomisation unit} is what is assigned: an order, a stock-day, a day, a client. \emph{Interference}\index{interference}
occurs when the outcome of one unit depends on the assignment of others, so that the difference between the arms of a test differs from the effect of
switching every unit. A \emph{sample ratio mismatch}\index{sample ratio mismatch} (SRM) is an observed split between the arms that differs, beyond chance,
from the planned one (Fabijan and co-authors).
\end{definition}

The chapter's simulated flow has 2\,400 parent orders a day in 500 stocks, drawn with popularity inversely proportional to rank. An order's \omterm{def:rs:simulation-versus-live:calibration}{implementation
shortfall}, in basis points against its \omterm{def:rs:simulation-versus-live:calibration}{arrival price}, is the firm's pre-trade cost estimate (standard deviation 6.9 across orders), the estimate's error,
the beta-adjusted market move over the order's scheduled window signed by its side, idiosyncratic noise, and a day effect common to all orders; the mean is
11.7 and the standard deviation 23.1. The new algorithm saves 0.30 basis point on each order it handles, and adds 0.25 to the other orders of the same stock,
side and day when all of them are on the new algorithm (it trades earlier in the window and leaves its siblings more impact).

\paragraph{Assignment.} \texttt{firm.abtest.assign} hashes the unit's identifier with the experiment's name and compares the first eight bytes of the SHA-256
digest, read as a fraction, with the planned share. The assignment is deterministic (the same order always goes to the same arm, on every server, without a
shared table), independent across experiments with different names, and reproducible afterwards by anyone holding the identifiers. In the two-week
experiment of the opening, 2\,485 of 24\,000 orders went to the new algorithm against a planned 2\,400: $z = 1.83$, a p-value of 0.07, no mismatch at 5\%. A
mismatch would not be a statistical curiosity: it is the symptom of orders lost, duplicated or re-routed in one arm, and it invalidates the comparison until
explained.

\paragraph{\omterm{def:rs:live-experiments:ab}{Interference}.} In the flow, 88\% of orders share their stock, side and day with another of the firm's orders. When orders are randomised one by
one, a treated order and a control order have, on average, the same share of treated siblings, so the difference between the arms measures the direct
saving, $-0.30$. The desk's question is what happens when every order moves: $-0.30 + 0.25 \times 0.883 = -0.08$. A \omterm{def:rs:live-experiments:ab}{randomisation unit} must contain the
\omterm{def:rs:live-experiments:ab}{interference}: here the stock-day, all of a stock's orders on a day in the same arm. The day itself contains every \omterm{def:rs:live-experiments:ab}{interference} within a day, but offers
only one unit a day, and the day effect is then the noise; alternating the treatment through time in this way is a switchback design (Bojinov, Simchi-Levi
and Zhao), common where units interfere, at the price of few units and of carryover from one period to the next.

\omcode{code/firm/abtest/firm_abtest.py}{74}{91}{The comparison when whole clusters are assigned: order-weighted means, cluster-robust standard error.}\label{lst:rs:live-experiments:cluster}

\section{Variance reduction}

\begin{definition}[CUPED]\label{def:rs:live-experiments:cuped}
\emph{CUPED}\index{CUPED} (controlled experiments using pre-experiment data; Deng, Xu, Kohavi and Walker) compares the arms on the adjusted outcome
$y - (x - \bar x)^\top\theta$, where $x$ are covariates the treatment cannot affect and $\theta$ is the coefficient of the regression of $y$ on $x$ pooled over
both arms.
\end{definition}

\begin{proposition}[What CUPED buys]\label{prop:rs:live-experiments:cuped}
The adjusted difference in means is unbiased for the treatment effect, and its variance is the unadjusted variance times $1 - R^2$, where $R^2$ is the share of
the outcome's variance explained by the covariates.
\end{proposition}

\begin{proof}
Randomisation gives the covariates the same distribution in both arms, so $\mathbb{E}[\bar x_T - \bar x_C] = 0$ and the adjustment removes nothing from the
expected difference. The variance of $y - x^\top\theta$ is smallest at the regression coefficient, where it is $\operatorname{Var}(y)(1 - R^2)$; the difference
of two independent means inherits the factor. Estimating $\theta$ on the pooled data changes this by a term of order $1/n$.
\end{proof}

The covariates must be fixed before assignment or lie outside the firm's reach. The pre-trade cost estimate qualifies, and explains 9\% of the variance of the
shortfall. So does the market's beta-adjusted move over the order's scheduled window: the portfolio manager fixes the window before the order is
assigned, and the market does not move because one of the firm's algorithms rather than the other is working it. With both covariates $R^2$ is 0.51, close
to the halving Deng and co-authors reported for Bing's experiments. The order's realised duration, its fill rate or its participation do not qualify: the
algorithm changes them, and adjusting for them removes part of the effect. Post-stratification (Miratrix, Sekhon and Yu) is \omterm{def:rs:live-experiments:cuped}{CUPED} with stratum indicators as
covariates, nearly as efficient as assigning within strata in advance; for a design randomised by stock-day, analysing within each day (subtracting each day's
mean from outcome and covariates) is the stratification that removes the day effect.

Each design below was run as 100 simulated experiments of 60 days on half the flow (the mean estimate, the spread of the estimates, the mean
reported standard error, and the days that standard error implies for 80\% power on 0.3 basis point):

\begin{center}
\small
\begin{tabular}{@{}lrrrr@{}}
\toprule
design and estimator & estimate & spread & s.e. & days for 0.3 bp\\
\midrule
by order: difference in means & $-0.30$ & 0.123 & 0.122 & 77\\
by order: stratified (pre-trade quintiles) & $-0.30$ & 0.119 & 0.117 & 72\\
by order: \omterm{def:rs:live-experiments:cuped}{CUPED}, pre-trade estimate & $-0.30$ & 0.121 & 0.116 & 70\\
by order: \omterm{def:rs:live-experiments:cuped}{CUPED}, pre-trade and market & $-0.29$ & 0.095 & 0.085 & 38\\
by day & $-0.16$ & 0.656 & 0.788 & 3\,253\\
by stock-day & $-0.09$ & 0.235 & 0.235 & 290\\
by stock-day, within day & $-0.09$ & 0.188 & 0.184 & 178\\
by stock-day, within day, \omterm{def:rs:live-experiments:cuped}{CUPED} & $-0.08$ & 0.093 & 0.082 & 35\\
\bottomrule
\end{tabular}
\end{center}

Every order-level design estimates the direct saving, $-0.30$; every design that assigns whole stock-days or days estimates the rollout's, $-0.08$ (the day
design's $-0.16$ is within its spread of 0.66). The reported standard errors match the spread of the estimates across experiments. \omterm{def:rs:live-experiments:cuped}{CUPED} with the market
covariate halves the days needed; the pre-trade estimate alone and stratification on it buy little, because the estimate explains little. Randomising by
day multiplies the days by 42. Randomising by stock-day and analysing within the day, with \omterm{def:rs:live-experiments:cuped}{CUPED}, is as precise as the best order-level design and
answers the right question; the answer is small, and detecting a saving of 0.08 basis point at 80\% power would take 503 trading days.

\section{Sizing and stopping an experiment}

\begin{definition}[Minimum detectable effect, sequential test]\label{def:rs:live-experiments:mde}
The \emph{minimum detectable effect}\index{minimum detectable effect} of a design is the smallest true effect it detects with a stated power at a stated size.
A \emph{sequential test}\index{sequential test} examines the data as they arrive and may stop at any look, with its size controlled over all the looks it
might take.
\end{definition}

\begin{proposition}[Minimum detectable effect and sample size]\label{prop:rs:live-experiments:mde}
For a two-sided test of size $\alpha$ and power $1-\beta$ on an estimate with standard error $\mathrm{se}$, the \omterm{def:rs:live-experiments:mde}{minimum detectable effect} is
$(z_{1-\alpha/2} + z_{1-\beta})\,\mathrm{se}$, which is $2.80\,\mathrm{se}$ at 5\% and 80\%. With outcome standard deviation $\sigma$ and a share $s$ of $n$
units treated, detecting an effect $\Delta$ needs
\[
n = \Bigl(\frac{(z_{1-\alpha/2} + z_{1-\beta})\,\sigma}{\Delta}\Bigr)^2\Bigl(\frac1s + \frac1{1-s}\Bigr).
\]
\end{proposition}

\begin{proof}
The test rejects when $|\hat\Delta|/\mathrm{se} > z_{1-\alpha/2}$; if $\hat\Delta \sim N(\Delta, \mathrm{se}^2)$ this has probability $1-\beta$ (neglecting the
other tail) when $\Delta/\mathrm{se} = z_{1-\alpha/2} + z_{1-\beta}$. Substitute $\mathrm{se}^2 = \sigma^2/(sn) + \sigma^2/((1-s)n)$.
\end{proof}

The standard deviation of 23.1 basis points and 2\,400 orders a day give 77.7 days at half the flow for 0.3 basis point, and 38.2 with \omterm{def:rs:live-experiments:cuped}{CUPED}'s residual
standard deviation of 16.2; the simulated experiments, rejecting in 78\% of runs at 80 days without \omterm{def:rs:live-experiments:cuped}{CUPED} and 76\% at 40 days with it, confirm the formula
(\cref{fig:rs:live-experiments:power}). The equal split minimises the sample size; a 10\% experiment, as in the opening, needs $(1/0.1 + 1/0.9)/4 = 2.8$
times as many orders.

\begin{omfigure}
\begin{tikzpicture}
\begin{axis}[width=11cm,height=5.6cm,xlabel={trading days at half the flow},ylabel={power for 0.3 bp},tick label style={font=\small},label style={font=\small},
  xmin=0,xmax=120,ymin=0,ymax=1.02,grid=major,grid style={gray!25},legend columns=3,
  legend style={font=\footnotesize,at={(0.5,-0.3)},anchor=north,draw=none,fill=none,/tikz/every even column/.append style={column sep=8pt}},
  table/col sep=comma]
\addplot[black!60,thick] table[x=days,y=raw]{figdata/research/21-live-experiments/power.csv};
\addplot[omDef,very thick] table[x=days,y=cuped]{figdata/research/21-live-experiments/power.csv};
\addplot[omThm,thick,dashed] table[x=days,y=day]{figdata/research/21-live-experiments/power.csv};
\addplot[only marks,mark=o,mark size=2pt,black!60,forget plot] table[x=days,y=raw]{figdata/research/21-live-experiments/power_sim.csv};
\addplot[only marks,mark=*,mark size=1.8pt,omDef,forget plot] table[x=days,y=cuped]{figdata/research/21-live-experiments/power_sim.csv};
\draw[black!50,densely dotted] (axis cs:0,0.8) -- (axis cs:120,0.8);
\legend{by order,by order with CUPED,by day}
\end{axis}
\end{tikzpicture}

\omcaption{Power to detect a saving of 0.3 basis point against the length of the experiment: the formula at each design's standard error (lines) and the share
of 100 simulated experiments that rejected (dots). The dotted line is 80\%. Data: \texttt{rs\_abtest}.}\label{fig:rs:live-experiments:power}
\end{omfigure}

\paragraph{Peeking.} A desk watching an experiment's dashboard every morning and stopping the first day its p-value falls below 5\% runs a different test
from the one the p-value describes. On 4\,000 simulated experiments with no effect at all, a fixed-horizon test looked at once rejects in 5\% of them; looked
at every day, it rejects at least once in 29\% within 38 days, 33\% within 60 and 38\% within 120 (\cref{fig:rs:live-experiments:peeking}).

The mixture sequential probability ratio test, a descendant of Wald's \omterm{def:rs:live-experiments:mde}{sequential test} (1945), turns continuous monitoring into a valid procedure (Johari,
Koomen, Pekelis and Walsh, who deployed it on a commercial testing platform). For an estimate $\hat\Delta$ with variance $V$ at the current look, and effects
drawn from $N(0, \tau^2)$ under the alternative, the mixture likelihood ratio against no effect is
\[
\Lambda = \sqrt{\frac{V}{V+\tau^2}}\,\exp\Bigl(\frac{\tau^2\hat\Delta^2}{2V(V+\tau^2)}\Bigr).
\]
Under no effect $\Lambda$ is a nonnegative martingale starting at 1, so the probability that it ever exceeds $1/\alpha$ is at most $\alpha$ (Ville's
inequality). The always-valid p-value is the running minimum of $1/\Lambda$; stopping whenever it falls below $\alpha$ keeps the error at $\alpha$ however often
one looks. On the same null experiments it crossed 5\% in 46, 62 and 82 of the 4\,000 by days 38, 60 and 120. With the real saving of 0.30 and \omterm{def:rs:live-experiments:cuped}{CUPED}, and
$\tau = 0.3$, half the experiments stop by day 38, the fixed horizon's length, 4\% by day 10 and 98\% by day 120. A harmful change of $+1$ basis point stops at
a median of 6 days. The \omterm{def:rs:live-experiments:mde}{sequential test} gives up a fixed end date; it keeps the size under monitoring and stops early when the effect is large, which is when
stopping matters.

\begin{omfigure}
\begin{tikzpicture}
\begin{axis}[name=a,width=6.8cm,height=5.4cm,xlabel={\small days of daily looks},ylabel={\small false-positive rate},tick label style={font=\footnotesize},
  xmin=0,xmax=120,ymin=0,ymax=0.42,grid=major,grid style={gray!25},title={\small no effect},
  yticklabel style={/pgf/number format/fixed,/pgf/number format/precision=2},scaled y ticks=false,table/col sep=comma]
\addplot[black!60,thick] table[x=days,y=naive]{figdata/research/21-live-experiments/peeking.csv};
\addplot[omDef,very thick] table[x=days,y=msprt]{figdata/research/21-live-experiments/peeking.csv};
\draw[omThm,densely dotted,thick] (axis cs:0,0.05) -- (axis cs:120,0.05);
\end{axis}
\begin{axis}[at={(a.east)},anchor=west,xshift=1.3cm,width=6.8cm,height=5.4cm,xlabel={\small day},ylabel={\small share stopped},
  tick label style={font=\footnotesize},xmin=0,xmax=120,ymin=0,ymax=1.02,grid=major,grid style={gray!25},title={\small saving of 0.30 bp},
  table/col sep=comma]
\addplot[omDef,very thick] table[x=days,y=share]{figdata/research/21-live-experiments/stopping.csv};
\draw[black!60,dashed] (axis cs:38,0) -- (axis cs:38,1.02);
\end{axis}
\end{tikzpicture}

\omcaption{Left: the share of 4\,000 null experiments declared significant at some daily look, with a fixed-horizon test (grey) and the mixture \omterm{def:rs:live-experiments:mde}{sequential
test} (blue); the dotted line is 5\%. Right: the day the \omterm{def:rs:live-experiments:mde}{sequential test} stops under the real saving, with \omterm{def:rs:live-experiments:cuped}{CUPED}; the dashed line is the 38-day fixed horizon
that has 80\% power. Data: \texttt{rs\_abtest.peeking}, \texttt{stopping}.}\label{fig:rs:live-experiments:peeking}
\end{omfigure}

\paragraph{The plan.} Everything that decides the outcome is written down before the first order is assigned: the metric and its guardrails (fill rate,
rejects, mark-outs), the \omterm{def:rs:live-experiments:ab}{randomisation unit}, the share, the covariates, the horizon or the sequential rule with its $\tau$, the \omterm{def:rs:live-experiments:mde}{minimum detectable effect},
and what will be done at each outcome. Chapter~1's \omterm{def:rs:the-research-process:log}{research log} records it; chapter~20's arithmetic about trials applies to experiments as much as to
\omterm{def:rs:vectorised-backtests:backtest}{backtests}, and a desk that runs twenty experiments a year will see one false saving a year at 5\%.

\section{Tutorial: is the new algorithm better?}\label{tut:rs:live-experiments:tut}

\begin{tutorial}
\textbf{Goal.} Simulate the firm's order flow with a new execution algorithm, assign orders by hash, compare designs and estimators, size the experiment, and
monitor it sequentially. \textbf{End state:} the table, \cref{fig:rs:live-experiments:power,fig:rs:live-experiments:peeking}.
\begin{tutsteps}
\item \textbf{Assignment and estimation}: the hash, the difference in means, \omterm{def:rs:live-experiments:cuped}{CUPED}.
\omcode{code/firm/abtest/firm_abtest.py}{35}{59}{Hash assignment, the difference in means and CUPED.}\label{lst:rs:live-experiments:assign}
\item \textbf{The flow}: \texttt{rs\_abtest.orders(days, seed, share, unit)} draws the orders, their covariates, the siblings and the outcome; \texttt{hook()}
runs the two-week, 10\% experiment and its SRM check.
\item \textbf{Designs}: \texttt{calibration()} runs 100 experiments per design; \texttt{rollout()} gives the true effect of switching everything.
\item \textbf{Sequential monitoring}: the mixture likelihood ratio and the always-valid p-value.
\omcode{code/firm/abtest/firm_abtest.py}{121}{129}{The mixture sequential probability ratio test.}\label{lst:rs:live-experiments:msprt}
\item \textbf{Run} \texttt{peeking()}, \texttt{stopping()} and \texttt{fig\_abtest.py}.
\end{tutsteps}
\textbf{What to change next.} Make the new algorithm harm its siblings more than it helps itself (a spillover of 0.40) and see an order-level test approve a
change that loses money; randomise by client instead of stock-day and ask what \omterm{def:rs:live-experiments:ab}{interference} that contains.
\end{tutorial}

\section{Build: the experiment toolkit}\label{bld:rs:live-experiments:abtest}

\begin{build}
\bfield{Purpose} Every change to a trading component reaches full flow through a \omterm{def:rs:live-experiments:canary}{staged rollout}, and every claimed improvement comes from a randomised
experiment, analysed as planned.
\bfield{Interface} \texttt{bucket(unit, experiment, salt)}, \texttt{assign(units, experiment, share, salt)}, \texttt{diff\_means(y, t)},
\texttt{cuped(y, X, t)}, \texttt{stratified(y, t, strata)}, \texttt{cluster\_diff(y, t, cluster)}, \texttt{demean\_by(v, group)},
\texttt{power(effect, se, alpha)}, \texttt{mde(se, alpha, power)}, \texttt{sample\_size(effect, sd, alpha, power, share)},
\texttt{srm\_pvalue(n\_t, n\_c, share)}, \texttt{msprt(d, se, tau)}, \texttt{always\_valid\_p(d, se, tau)}.
\bfield{Rules} Assignment by hash of a stable identifier, never by a mutable table; covariates fixed before assignment; the \omterm{def:rs:live-experiments:ab}{randomisation unit} contains the
\omterm{def:rs:live-experiments:ab}{interference}; the SRM check runs before any estimate is read; a fixed horizon is not looked at early, and a sequential rule is chosen before the start.
\bfield{Acceptance tests} \texttt{code/firm/abtest/tests/}: a deterministic, balanced assignment, independent across experiments; \omterm{def:rs:live-experiments:cuped}{CUPED}'s variance falling by
$1-R^2$ on a known covariate, stratification between the raw and adjusted estimates; cluster arms constant or an error; power, \omterm{def:rs:live-experiments:mde}{minimum detectable effect} and
sample size consistent; an always-valid p-value non-increasing and below 5\% in fewer than 5\% of 2\,000 null streams of 200 looks, while daily looks at
a fixed-horizon test exceed 30\%.
\bfield{Stretch} Variance-weighted estimators for unequal clusters; guardrail metrics with their own \omterm{def:rs:live-experiments:mde}{sequential tests}; a switchback design with carryover.
\end{build}

\begin{omsources}
\item A.~Deng, Y.~Xu, R.~Kohavi and T.~Walker, ``Improving the sensitivity of online controlled experiments by utilizing pre-experiment data'', \emph{WSDM}, 2013.
\item R.~Kohavi, D.~Tang and Y.~Xu, \emph{Trustworthy Online Controlled Experiments}, Cambridge University Press, 2020.
\item A.~Fabijan et al., ``Diagnosing sample ratio mismatch in online controlled experiments'', \emph{KDD}, 2019.
\item R.~Johari, P.~Koomen, L.~Pekelis and D.~Walsh, ``Always valid inference: continuous monitoring of A/B tests'', \emph{Operations Research} 70(3), 2022
(and ``Peeking at A/B tests'', \emph{KDD}, 2017).
\item A.~Wald, ``Sequential tests of statistical hypotheses'', \emph{Annals of Mathematical Statistics} 16(2), 1945.
\item SEC, \emph{In the Matter of Knight Capital Americas LLC}, Release No.~34-70694, 16 October 2013.
\item I.~Bojinov, D.~Simchi-Levi and J.~Zhao, ``Design and analysis of switchback experiments'', \emph{Management Science} 69(7), 2023.
\item L.~W.~Miratrix, J.~S.~Sekhon and B.~Yu, ``Adjusting treatment effect estimates by post-stratification in randomized experiments'', \emph{Journal of the
Royal Statistical Society B} 75(2), 2013.
\end{omsources}

\section{Exercises}

\begin{exercise}[$\star$]\label{exo:rs:live-experiments:1}
The two-week experiment estimates a saving of 0.55 basis point with a standard error of 0.48. Give the 95\% interval and the verdict.
\end{exercise}

\begin{exercise}[$\star$]\label{exo:rs:live-experiments:2}
The same experiment planned 10\% of 24\,000 orders and assigned 2\,485. Compute the SRM statistic and its p-value. What would you do with a p-value of 0.0001?
\end{exercise}

\begin{exercise}[$\star$]\label{exo:rs:live-experiments:3}
An estimate's standard error is 0.085 basis point. What is the \omterm{def:rs:live-experiments:mde}{minimum detectable effect} at 5\% and 80\%, and what is the power against a true saving of 0.3?
\end{exercise}

\begin{exercise}[$\star\star$]\label{exo:rs:live-experiments:4}
Which of these may serve as \omterm{def:rs:live-experiments:cuped}{CUPED} covariates in an execution experiment: the pre-trade cost estimate; the market's move over the order's scheduled window; the
order's realised duration; the algorithm's fill rate; the stock's shortfall on the previous day? With covariates explaining 51\% of the variance, how many of
77 days remain?
\end{exercise}

\begin{exercise}[$\star\star$]\label{exo:rs:live-experiments:5}
A defect costs 5 basis points an order, the per-order standard deviation is 23.1, the alarm is at $z = 3$ and the flow is 2\,400 orders a day. How many
orders and how long until the alarm with a 1\% canary, a 10\% one and a 50\% one? What would have stopped Knight's defect, and why is it not statistics?
\end{exercise}

\begin{exercise}[$\star\star$]\label{exo:rs:live-experiments:6}
With a direct saving of 0.30, a spillover of 0.25 onto fully treated siblings and 88.3\% of orders with siblings, what is the effect of the full rollout? Why
does randomising by order measure 0.30, and which unit measures the rollout?
\end{exercise}

\begin{exercise}[$\star\star\star$]\label{exo:rs:live-experiments:7}
\emph{Coding.} Run \texttt{rs\_abtest.stopping} under the real saving with $\tau = 0.1$ and $\tau = 1$ as well as 0.3. Report the median stopping days and
explain why both a small and a large $\tau$ are slower.
\end{exercise}

\begin{exercise}[$\star\star\star$]\label{exo:rs:live-experiments:8}
\emph{Find the flaw.} ``We checked the experiment's dashboard every morning, and on day 23 the p-value fell to 0.04, so we stopped and rolled the new
algorithm out: a significant saving at 5\%.''
\end{exercise}

\section{Problem: Is the New Algorithm Better?}

\begin{problem}[{Weekend problem --- an execution experiment, designed and read}]\label{pb:rs:live-experiments:1}
The chapter's simulated flow: 2\,400 orders a day, a new algorithm saving 0.30 basis point on its own orders and costing its siblings 0.25.

\medskip\noindent\textbf{Part I --- The design.}
\begin{enumerate}
\item What are the mean and standard deviation of an order's shortfall, and what are its components?
\item What properties does hash assignment give, and why do they matter on a trading system?
\item Check the opening experiment's split.
\item What does randomising by order measure, and what does the full rollout save?
\item What should Knight's staged deployment have checked at each stage?
\end{enumerate}

\medskip\noindent\textbf{Part II --- Precision.}
\begin{enumerate}[resume]
\item How many days at half the flow does the difference in means need for 0.3 basis point?
\item Which covariates does \omterm{def:rs:live-experiments:cuped}{CUPED} use, what is their $R^2$, and how many days remain?
\item What do the pre-trade estimate and stratification on it buy, and why so little?
\item What does randomising by day cost?
\item How precise is the stock-day design analysed within the day with \omterm{def:rs:live-experiments:cuped}{CUPED}, and what does it estimate?
\end{enumerate}

\medskip\noindent\textbf{Part III --- Stopping.}
\begin{enumerate}[resume]
\item What false-positive rates does daily peeking at a fixed-horizon test give by days 38 and 60?
\item How many of the 4\,000 null experiments does the mixture \omterm{def:rs:live-experiments:mde}{sequential test} stop by the same days?
\item Under the real saving, when does the \omterm{def:rs:live-experiments:mde}{sequential test} stop?
\item How fast does it catch a harmful change of 1 basis point?
\end{enumerate}

\medskip\noindent\textbf{Part IV --- The verdict.}
\begin{enumerate}[resume]
\item State the \emph{named result}: the trading days needed to detect a 0.3 basis point improvement at 80\% power, with and without \omterm{def:rs:live-experiments:cuped}{CUPED}.
\item Is the new algorithm better?
\item How would you roll it out?
\item How long would a 1\% canary take to notice a defect costing 5 basis points an order, and what would you rely on instead?
\item What goes into the experiment's plan before the first order?
\item In one sentence: what does an \omterm{def:rs:live-experiments:ab}{A/B test} on live flow measure?
\end{enumerate}
\end{problem}

\section{Interview questions}

\begin{interviewq}[$\star$ \iqroles{trader, developer}]\label{iq:rs:live-experiments:1}
How would you put a new execution algorithm into production safely?
\end{interviewq}

\begin{interviewq}[$\star\star$ \iqroles{researcher}]\label{iq:rs:live-experiments:2}
What is \omterm{def:rs:live-experiments:cuped}{CUPED}, when does it help, and what can go wrong with it?
\end{interviewq}

\begin{interviewq}[$\star\star$ \iqroles{researcher}]\label{iq:rs:live-experiments:3}
Why is it a problem to check an \omterm{def:rs:live-experiments:ab}{A/B test} every day and stop when it is significant, and what do you do instead?
\end{interviewq}

\begin{interviewq}[$\star\star$ \iqroles{researcher, trader}]\label{iq:rs:live-experiments:4}
An \omterm{def:rs:live-experiments:ab}{A/B test} randomised by order showed the new algorithm saving 0.3 basis point. After the full rollout, costs did not fall. What happened?
\end{interviewq}

\begin{interviewq}[$\star\star$ \iqroles{developer}]\label{iq:rs:live-experiments:5}
How do you assign orders to the arms of an experiment across many servers, and how do you check that the assignment worked?
\end{interviewq}

\begin{interviewq}[$\star\star\star$ \iqroles{researcher}]\label{iq:rs:live-experiments:6}
Derive the sample size needed to detect an effect $\Delta$ with power $1-\beta$. With a standard deviation of 23.1 basis points and 2\,400 orders a day, how
many days for 0.3 basis point?
\end{interviewq}
